% GENERATED FILE. Edit canonical lessons, then run npm run generate:books.
% canonical-source-hash: fnv1a64:ad5a733afa7ca6d7
% canonical-lessons: AR-C36-zayt, AR-C36-kub, AR-C36-tabaq, AR-C36-milaqa

\chapter{Cups, Plates, and More Than One}
\label{ch:ar-cups-plates-and-more-than-one}
\hlchaptermodality{36}

\begin{chapteropening}
\textbf{By the end of this chapter:} \emph{I can ask for oil, a cup, a plate and a spoon in Arabic, and I can tell a broken plural from a sound one and say which kind of noun usually takes which.}

Set the mi- tool shape of milʿaqa and mirʾāh against the ma- place shape of maṭʿam and maktab, then sort the plurals: akwāb, aṭbāq and malāʿiq re-poured into a new shape against qahawāt and ṣadīqāt, which only add an ending — and ask for each thing with min faḍlik.
\end{chapteropening}

\section[zayt]{\ar{زيت} (zayt) — \textquotedblleft{}oil,\textquotedblright{} and a letter you already knew}

\label{lesson:AR-C36-zayt}



\begin{quote}
One new letter to name — and you will find you already own it, because it is an old letter wearing a dot.
\end{quote}

\subsection*{The letters in this word}
\textbf{\ar{ز}} is \textbf{zāy}, the \emph{z}. And here is all there is to it: \textbf{\ar{ز}} is \textbf{\ar{ر}} (\emph{rā}) with \textbf{one dot above}.

That is the abjad's favourite trick. \textbf{\ar{ب}}, \textbf{\ar{ن}} and \textbf{\ar{ت}} share one bowl and differ by dots. \textbf{\ar{ح}}, \textbf{\ar{خ}} and \textbf{\ar{ج}} share one hook. Now \textbf{\ar{ر}} and \textbf{\ar{ز}} share one curve.

Everything \emph{rā} does, \emph{zāy} does: it hangs below the line and \textbf{refuses to join} what follows — the non-joiner that ended \textbf{\ar{سكر}}.

Then \textbf{\ar{ي}} (\emph{yāʾ}) and \textbf{\ar{ت}} (\emph{tāʾ}). So \textbf{\ar{زيت}} is \textbf{\ar{ز}} alone, then \textbf{\ar{يت}} joined.

\begin{cousinweb}
\textbf{\ar{زَيْت}} (\emph{zayt}) = \textquotedblleft{}\textbf{oil}\textquotedblright{} — in the first place olive oil, and by extension oil of any kind. Root \textbf{\ar{ز}-\ar{ي}-\ar{ت}}, which also gives:

\begin{itemize}
\raggedright
  \item \textbf{\ar{زَيْتون}} (\emph{zaytūn}) — \textquotedblleft{}\textbf{olives}.\textquotedblright{}
  \item \textbf{\ar{زَيْتونة}} (\emph{zaytūna}) — \textquotedblleft{}\textbf{an olive},\textquotedblright{} one of them, with the tied \textbf{\ar{ة}} making the single thing out of the collective.
\end{itemize}

And now Spain. You met \textbf{azúcar} and found Arabic's \textbf{al-} welded onto its front. The same thing happened here, twice over:

\begin{itemize}
\raggedright
  \item \textbf{\ar{الزَّيت}} (\emph{az-zayt}) became Spanish \textbf{aceite}, \textquotedblleft{}oil.\textquotedblright{}
  \item \textbf{\ar{الزَّيتونة}} (\emph{az-zaytūna}) became Spanish \textbf{aceituna}, \textquotedblleft{}an olive.\textquotedblright{}
\end{itemize}

\textbf{\ar{ز}} is a \textbf{sun letter}, which is why \emph{al-} is said \emph{az-} — the assimilation that makes \emph{al-} + \emph{sukkar} into \emph{as-sukkar}. Spanish took the sound it heard, article and all. Portuguese did the same with \textbf{azeite}.
\end{cousinweb}

\subsection*{Your turn}
Say these aloud:

\begin{itemize}
\raggedright
  \item zāy is rā plus one dot above, and it breaks as rā breaks
  \item the dot families — bāʾ, nūn, tāʾ; ḥāʾ, khāʾ, jīm; rā and zāy
  \item zāy, yāʾ, tāʾ — \textquotedblleft{}zayt,\textquotedblright{} oil
  \item the olives — zaytūn; one olive, zaytūna
  \item the welded article — az-zayt gave aceite, as-sukkar gave azúcar
  \item the friend and his root — ṣadīq, from ṣ-d-q, to be true
  \item the pair one letter apart — ibn, bint
\end{itemize}

\begin{tcolorbox}[breakable,colback=teal!4,colframe=teal!35!black,title={Before you move on}]
What is \textbf{\ar{ز}}, described from a letter you had? (\emph{\textbf{Rā}} \textbf{with one dot above} — and it breaks the same way.) What does \textbf{\ar{ز}-\ar{ي}-\ar{ت}} give besides oil? (\emph{\textbf{Zaytūn}} and \emph{\textbf{zaytūna}}.) Which Spanish words carry this root with Arabic's article stuck on, and which one did \emph{sukkar} give? (\emph{\textbf{Aceite}} and \emph{\textbf{aceituna}}; \emph{\textbf{azúcar}}.) Why \emph{az-} rather than \emph{al-}? (\textbf{\ar{ز} is a sun letter}.)
\end{tcolorbox}
\section[kūb]{\ar{كوب} (kūb) — \textquotedblleft{}a cup,\textquotedblright{} and more than one of them}

\label{lesson:AR-C36-kub}



\begin{quote}
A short, useful word — and the door into the thing about Arabic nouns that surprises every learner.
\end{quote}

\subsection*{The letters in this word}
Nothing new. \textbf{\ar{ك}} (\emph{kāf}), the angular \emph{k}; \textbf{\ar{و}} (\emph{wāw}), here carrying the long \textbf{ū}; and \textbf{\ar{ب}} (\emph{bāʾ}).

\textbf{\ar{و}} is a non-joiner, so \textbf{\ar{كوب}} goes \textbf{\ar{كو}}, break, \textbf{\ar{ب}}. Another non-joiner: \textbf{\ar{ر}} ended \emph{sukkar}, \textbf{\ar{ز}} opened \emph{zayt}, and \textbf{\ar{و}} splits this.

\begin{cousinweb}
\textbf{\ar{كوب}} (\emph{kūb}) = \textquotedblleft{}\textbf{a cup}.\textquotedblright{} Root \textbf{\ar{ك}-\ar{و}-\ar{ب}}.

Now say \textquotedblleft{}cups.\textquotedblright{} In English you add an \emph{s}. In Arabic you do something else entirely:

\textbf{\ar{أَكْواب}} (\emph{akwāb}).

Nothing was added to the end. The root \textbf{\ar{ك}-\ar{و}-\ar{ب}} was taken out of one shape and \textbf{poured into another}, \textbf{\ar{أَفْعال}} (\emph{afʿāl}). Arabic calls this a \textbf{broken plural}, and it is not an exception — it is how most ordinary Arabic nouns make their plurals.

You have seen a root move between shapes many times: \textbf{\ar{ك}-\ar{ت}-\ar{ب}} became \emph{kitāb}, \emph{kātib}, \emph{maktab}, \emph{maktaba}. A broken plural is that same machinery pointed at number instead of meaning. So this is not a new rule to learn. It is a rule you already have, doing a job you did not expect it to do.

On where \emph{kūb} came from, honesty first. It sits close to a cup-word that travelled all round the Mediterranean — Latin \textbf{cuppa}, behind English \textbf{cup} and French \textbf{coupe}. Whether Arabic took it from that family, or the resemblance is coincidence, is \textbf{not settled}: a resemblance worth knowing, not a descent worth claiming.
\end{cousinweb}

\subsection*{Your turn}
Say these aloud:

\begin{itemize}
\raggedright
  \item kāf, wāw, bāʾ — \textquotedblleft{}kūb,\textquotedblright{} a cup
  \item more than one — \textquotedblleft{}akwāb,\textquotedblright{} and nothing was added to the end
  \item what happened instead — the root was poured into the shape afʿāl
  \item the same machinery you had — kitāb, kātib, maktab, maktaba
  \item ask for one — \textquotedblleft{}kūb māʾ, min faḍlik\textquotedblright{} — a cup of water, from your grace
  \item and answer the thanks — \textquotedblleft{}ʿafwan\textquotedblright{}
  \item the odd old word for water — māʾ, whose plural miyāh follows no pattern
  \item the non-joiners — zāy in zayt, dāl in ṣadīq, wāw here
\end{itemize}

\begin{tcolorbox}[breakable,colback=teal!4,colframe=teal!35!black,title={Before you move on}]
How does Arabic make \textquotedblleft{}cups\textquotedblright{} out of \textbf{\ar{كوب}}? (\emph{\textbf{Akwāb}} — the root is \textbf{re-poured into a new shape}, not given an ending. That is a \textbf{broken plural}.) Which earlier root shows the same shape-changing machinery? (\textbf{\ar{ك}-\ar{ت}-\ar{ب}} — \emph{kitāb}, \emph{kātib}, \emph{maktab}, \emph{maktaba}.) Ask for a cup of water politely. (\emph{\textbf{Kūb māʾ, min faḍlik}}.) And reply to thanks? (\emph{\textbf{ʿAfwan}}.) Did Arabic get \emph{kūb} from Latin \emph{cuppa}? (\textbf{Unsettled} — a resemblance, not a claim.)
\end{tcolorbox}
\section[ṭabaq]{\ar{طبق} (ṭabaq) — \textquotedblleft{}a plate,\textquotedblright{} and a broken plural that repeats}

\label{lesson:AR-C36-tabaq}



\begin{quote}
One broken plural is a surprise. Two in the same shape is a pattern — and a pattern you can use.
\end{quote}

\subsection*{The letters in this word}
Nothing new. \textbf{\ar{ط}} (\emph{ṭāʾ}) is the emphatic \textbf{\ar{ت}}, named when you built \textbf{\ar{طعام}} (\emph{ṭaʿām}) — the dark, flat-tongued \emph{t}, standing beside \textbf{\ar{ص}} under \textbf{\ar{س}} and \textbf{\ar{ض}} under \textbf{\ar{د}}.

Then \textbf{\ar{ب}} (\emph{bāʾ}) and \textbf{\ar{ق}} (\emph{qāf}). All three join, so \textbf{\ar{طبق}} runs unbroken — unlike \textbf{\ar{كوب}}, which had to stop at its \emph{wāw}.

\begin{cousinweb}
\textbf{\ar{طَبَق}} (\emph{ṭabaq}) = \textquotedblleft{}\textbf{a plate, a dish}.\textquotedblright{} Root \textbf{\ar{ط}-\ar{ب}-\ar{ق}}.

Its plural is \textbf{\ar{أَطْباق}} (\emph{aṭbāq}).

Say that beside \textbf{\ar{أَكْواب}} (\emph{akwāb}) and listen. Same opening \textbf{\ar{أَ}}, same long \textbf{ā} before the last letter, same nothing-added-at-the-end. Both are the shape \textbf{\ar{أَفْعال}} (\emph{afʿāl}). So the broken plural is not chaos: it is a small set of shapes, and you have now met the commonest one twice.

The root itself means \textbf{to cover, to lie in layers} — a plate being the flat thing that covers. Which is why the same root gives:

\begin{itemize}
\raggedright
  \item \textbf{\ar{طابِق}} (\emph{ṭābiq}) — \textquotedblleft{}\textbf{a storey}\textquotedblright{} of a building. A layer of house. And notice its shape: that is the \textbf{doer} form of \emph{ṣādiq} and \emph{kātib}.
  \item \textbf{\ar{طَبَقة}} (\emph{ṭabaqa}) — \textquotedblleft{}\textbf{a layer},\textquotedblright{} and so \textquotedblleft{}\textbf{a social class},\textquotedblright{} with the tied \textbf{\ar{ة}} on the end.
\end{itemize}

A plate, a floor and a social class, off one idea about layers. And the plates themselves belong somewhere you already built: \textbf{\ar{مَطْعَم}} (\emph{maṭʿam}), the place of eating.
\end{cousinweb}

\subsection*{Your turn}
Say these aloud:

\begin{itemize}
\raggedright
  \item ṭāʾ, bāʾ, qāf — \textquotedblleft{}ṭabaq,\textquotedblright{} a plate
  \item the emphatic set once more — ṣād, ḍād, ṭāʾ
  \item the two broken plurals — akwāb, aṭbāq, both in the shape afʿāl
  \item what the root means — to cover, to lie in layers
  \item its other children — ṭābiq, a storey; ṭabaqa, a layer or a class
  \item where the plates live — maṭʿam, the place of eating
  \item the oil, whose zāy is rā plus a dot — zayt, zaytūn
\end{itemize}

\begin{tcolorbox}[breakable,colback=teal!4,colframe=teal!35!black,title={Before you move on}]
What shape do \textbf{\ar{أكواب}} and \textbf{\ar{أطباق}} share? (\emph{\textbf{Afʿāl}} — the broken plural you have now met twice.) What does \textbf{\ar{ط}-\ar{ب}-\ar{ق}} mean? (\textbf{To cover, to lie in layers}.) Name two more words it builds. (\emph{\textbf{Ṭābiq}}, a storey, and \emph{\textbf{ṭabaqa}}, a layer or a class.) Which plain letter is \textbf{\ar{ط}} the emphatic partner of? (\emph{\textbf{Tāʾ}}.) Where would you find a stack of \emph{aṭbāq}? (\textbf{In a \emph{maṭʿam}} — the place of eating.)
\end{tcolorbox}
\section[milʿaqa]{\ar{ملعقة} (milʿaqa) — \textquotedblleft{}a spoon,\textquotedblright{} the tool shape, and two ways to say more than one}

\label{lesson:AR-C36-milaqa}



\begin{quote}
The last word, and it pays two debts at once: the shape that makes tools, and the question of how Arabic counts past one.
\end{quote}

\subsection*{The letters in this word}
Nothing new: \textbf{\ar{م}} (\emph{mīm}), \textbf{\ar{ل}} (\emph{lām}), \textbf{\ar{ع}} (\emph{ʿayn}, the throat sound), \textbf{\ar{ق}} (\emph{qāf}) — and then \textbf{\ar{ة}}, the tied \emph{tāʾ}.

That \textbf{\ar{ة}} tells you, before anything else, that \textbf{\ar{ملعقة}} is \textbf{feminine} — the signal you read off \textbf{\ar{قهوة}} (\emph{qahwa}), \textbf{\ar{صديقة}} (\emph{ṣadīqa}) and \textbf{\ar{طَبَقة}} (\emph{ṭabaqa}).

\begin{cousinweb}
\textbf{\ar{مِلْعَقة}} (\emph{milʿaqa}) = \textquotedblleft{}\textbf{a spoon}.\textquotedblright{} Root \textbf{\ar{ل}-\ar{ع}-\ar{ق}}, whose verb \textbf{\ar{لَعِقَ}} (\emph{laʿiqa}) means \textquotedblleft{}\textbf{he licked, he lapped}.\textquotedblright{} A spoon is the thing you lap with.

Two shapes now stand side by side, differing by one vowel:

\begin{itemize}
\raggedright
  \item \textbf{\ar{مَـ}}, the \textbf{place} shape: \textbf{\ar{مَكْتَب}} (\emph{maktab}), the place of writing; \textbf{\ar{مَطْعَم}} (\emph{maṭʿam}), the place of eating.
  \item \textbf{\ar{مِـ}}, the \textbf{tool} shape: \textbf{\ar{مِرْآة}} (\emph{mirʾāh}), the thing you see with; \textbf{\ar{مِلْعَقة}} (\emph{milʿaqa}), the thing you lap with.
\end{itemize}

\emph{Ma-} builds a room. \emph{Mi-} puts something in your hand.

Now the plurals this chapter has been circling.

\textbf{Broken} — the root re-poured, nothing added: \textbf{\ar{أَكْواب}} (\emph{akwāb}), \textbf{\ar{أَطْباق}} (\emph{aṭbāq}), \textbf{\ar{مَلاعِق}} (\emph{malāʿiq}), spoons.

\textbf{Sound} — an ending added, the word left intact. Arabic uses it for the \textbf{\ar{ة}} words: \textbf{\ar{قَهْوة}} gives \textbf{\ar{قَهَوات}} (\emph{qahawāt}), and \textbf{\ar{صَديقة}} gives \textbf{\ar{صَديقات}} (\emph{ṣadīqāt}).

The rule of thumb: a noun in \textbf{\ar{ة}} usually takes the added ending, and most other everyday nouns break. \emph{Milʿaqa} is the honest reminder that \textquotedblleft{}usually\textquotedblright{} is doing work — it ends in \textbf{\ar{ة}} and breaks anyway.
\end{cousinweb}

\subsection*{Your turn}
Say these aloud:

\begin{itemize}
\raggedright
  \item mīm, lām, ʿayn, qāf, tāʾ marbūṭa — \textquotedblleft{}milʿaqa,\textquotedblright{} a spoon
  \item the verb under it — \textquotedblleft{}laʿiqa,\textquotedblright{} he lapped
  \item the place shape — maktab, maṭʿam
  \item the tool shape — mirʾāh, milʿaqa
  \item the broken plurals — kūb with its breaking wāw gives akwāb; aṭbāq; malāʿiq
  \item the sound plurals — qahwa gives qahawāt, ṣadīqa gives ṣadīqāt
  \item the tied tāʾ marking feminine — qahwa, ṣadīqa, ṭabaqa, milʿaqa
  \item ask for each — \textquotedblleft{}kūb, min faḍlik\textquotedblright{}; \textquotedblleft{}ṭabaq, min faḍlik\textquotedblright{}; \textquotedblleft{}milʿaqa, min faḍlik\textquotedblright{}
  \item and the oil — zayt, whose Spanish child is aceite
\end{itemize}

\begin{tcolorbox}[breakable,colback=teal!4,colframe=teal!35!black,title={Before you move on}]
What separates \textbf{\ar{مَـ}} from \textbf{\ar{مِـ}}? (\emph{\textbf{Ma-}} \textbf{makes a place}, \emph{\textbf{mi-}} \textbf{a tool} — \emph{maktab} and \emph{maṭʿam} against \emph{mirʾāh} and \emph{milʿaqa}.) What does \textbf{\ar{ل}-\ar{ع}-\ar{ق}} mean? (\textbf{To lap}.) Name this chapter's broken plurals, and a sound one. (\emph{\textbf{Akwāb}}, \emph{\textbf{aṭbāq}}, \emph{\textbf{malāʿiq}}; \emph{\textbf{qahawāt}} or \emph{\textbf{ṣadīqāt}}.) What does the \textbf{\ar{ة}} on \emph{milʿaqa} tell you? (\textbf{Feminine} — and unusually for a \textbf{\ar{ة}} word, it still breaks.) Ask for a spoon politely. (\emph{\textbf{Milʿaqa, min faḍlik}}.)
\end{tcolorbox}
